\scanpage{255}
\mspage{201}

on a un hom can.
\[
\begin{aligned}
\underline{\mathrm{hot}_{S,W}(X\times_S Y)}&\longrightarrow
\underline{\mathrm{hot}_{S,W}(X)\times_S\mathrm{hot}_{S,W}(Y)},\\
\Vert&\hspace{7em}\Vert\\
\mathrm{mult}_{X/S}(Y)&\qquad\mathrm{hot}_{S,W}(\unsure{\Phi|\mathscr{D}}(X)\times_S Y)
\end{aligned}
\]
induit par
\[
X\times_S Y\longrightarrow\Phi(X)\times_S Y
\]
mais cette\deleted{\unsure{ flèche}} dernière flèche n'a aucune raison d'être dans $W_S^d$, \deleted{\unsure{donc}} et n'induit \inserted{donc} pas \unsure{néc.$^{t}$} un iso dans $\mathrm{HOT}_W(S)$. (Contrexemple page 154.)

Il faut regarder à part le cas \uline{où $X$ est dans $\underline{\mathrm{Fib}}_W S$}, j'ai envie de dire que dans ce cas on a
\[
\begin{aligned}
\mathrm{hot}_{S,W}(X\times_S Y)&\overset{\sim}{\longrightarrow}
\underline{\mathrm{hot}_{S,W}(X)\times_S\mathrm{hot}_{S,W}(Y)}^{*}\\
&\hspace{5em}\Vert\\
&\quad\mathrm{hot}_{S,W}(\Phi(X)\times_S\Phi(Y))
\end{aligned}
\]
% slanted
\margin{Cette formule n'est établie qu'en supposant ou bien $Y$ $W$-lisse (en plus de $X$ $W$-fibré sur $S$ !), ou bien $X$ $W$-propre, p. ex. $X$ $W$-parfait \unsure{?|!}.}
pour \uline{tout} $Y$ dans $\mathrm{Cat}_{/S}$ \inserted{(pas néc. $W$-lisse)}, a fortiori que dans ce cas
\[
\left\{\begin{aligned}
\mathrm{mult}_{X/S}&=\mathrm{mult}_{\unsure{\xi|\zeta}}\qquad\text{dans }\mathrm{HOT}_W(S)\\
\text{où }\xi&=\mathrm{hot}_{S,W}(X/S),
\end{aligned}\right.\quad\bigl(\text{si }X\text{ dans }\underline{\mathrm{Fib}}_W\unsure{S|S'}\bigr)\ \Vert\tag*{$(**)$}
\]
\deleted{\unsure{mais il faut prouver que} \illegible{1}} \deleted{Quand} \inserted{Dans le cas où} $Y$ est lisse sur $S$ (ce qui suffit pour établir $(**)$), il suffit de \deleted{\unsure{voir}} \inserted{se rappeler} que $X\to\Phi(X)$, \deleted{étant dans} \deleted{$W_S^{\unsure{f}}$} dans $W_S^f$ (résulte du fait que $X$ est $W$-fibré sur $S^{*}$, donc $X_s\unsure{\mathbin{-}|\to}s\backslash X$ est dans $W$, $\forall s\in S$),
% slanted
\margin{il suffit $=$ $W$-prélisse sur $S$}
donc $\underbrace{X\times_S Y}_{X_Y}\to\underbrace{\Phi(X)}_{\Phi(X)_Y}\times_S Y$ est aussi dans $W_Y^f$, \deleted{et comme les deux sont $W$-prélisses sur $S$,} donc (comme \deleted{\unsure{$X,Y\in$}} $X_Y,\Phi(X)_Y$ sont $W$-fibrés sur $Y$) dans $W_Y^d$, donc aussi dans $W_S^d$. En fait, $=$ sans supposer $Y$ lisse, ceci montre que l'on a
% archive sheet number: 254
